A low surface brightness galaxy (LSB) is a diffuse galaxy with a surface
brightness that, when viewed from the Earth, is at least one magnitude lower
than the ambient night sky.

Some of its matter is in the form of ``baryonic'' matter such as neutral
hydrogen gas and stars. However, most of its matter is in the form of
invisible mass -- so called ``dark matter''. In this question, we will
investigate the mass of dark matter in a galaxy, the effect of dark matter on
the rotation curves of the galaxy, and the distribution of dark matter in the
galaxy.

The table below provides the data of a LSB galaxy named UGC4325. The galaxy is
assumed to be edge-on. At every distance $r$ from the centre of the galaxy, we
measure
\begin{enumerate}
\item $\lambda_{\mathrm{obs}}$, the observed wavelength of the H$\alpha$
  emission line. The Hubble expansion of the Universe has already been
  excluded from the data.
\item $V_{\mathrm{gas}}$, the contribution of the gas component to the
  rotation due to $M_{\mathrm{gas}}$, derived from HI surface densities.
\item $V_*$, the contribution of the stellar component to the rotation due to
  $M_*$, derived from R-band photometry.
\end{enumerate}
The rotational velocities of the test particle due to the gas component,
$V_{\mathrm{gas}}$, and the star component, $V_*$, are defined as the
velocities in the plane of the galaxy that would result from the corresponding
components without any external influences. These velocities are calculated
from the observed baryonic mass density distributions.

\begin{center}
\begin{tabular}{cccc}
$r$ & $\lambda_{\mathrm{obs}}$ & $V_{\mathrm{gas}}$ & $V_*$ \\
(kpc) & (nm) & (km/s) & (km/s) \\
\hline
 0.70 & 656.371 &  2.87 & 20.97 \\
 1.40 & 656.431 &  6.75 & 32.22 \\
 2.09 & 656.464 & 14.14 & 40.91 \\
 2.79 & 656.475 & 20.18 & 46.75 \\
 3.49 & 656.478 & 24.08 & 50.10 \\
 4.89 & 656.484 & 28.08 & 47.94 \\
 6.25 & 656.481 & 29.25 & 45.47 \\
 7.10 & 656.481 & 27.03 & 47.78 \\
 9.03 & 656.482 & 25.90 & 45.32 \\
12.05 & 656.482 & 21.03 & 42.30 \\
\end{tabular}
\end{center}

The mass of dark matter $M_{\mathrm{DM}}(r)$ within a volume of radius $r$ can
be defined in terms of the rotational velocity due to dark matter
$V_{\mathrm{DM}}$, the radius $r$ and the gravitational constant $G$,
\begin{equation}
M_{\mathrm{DM}}(r) = \frac{rV_{\mathrm{DM}}^2}{G}. \tag{1}
\end{equation}
To a good approximation, the observed rotational velocity $V_{\mathrm{obs}}$
can be modelled as
\begin{equation}
V_{\mathrm{obs}}^2 = V_{\mathrm{gas}}^2 + V_*^2 + V_{\mathrm{DM}}^2.
\tag{2}
\end{equation}
The observed rotational velocity $V_{\mathrm{obs}}$ depends on the mass of the
galaxy $M(r)$ within a volume of radius $r$ measured from the galaxy's centre.

The mass density $\rho_{\mathrm{DM}}(r)$ of dark matter within a volume of
radius $r$ can be modelled by a galaxy density profile,
\begin{equation}
\rho_{\mathrm{DM}}(r)
  = \frac{\rho_0}{1 + \left(\dfrac{r}{r_{\mathrm{C}}}\right)^2}
\tag{3}
\end{equation}
where $\rho_0$ and $r_{\mathrm{C}}$ are the central density and the core
radius of the galaxy, respectively. According to the density profile, the mass
of dark matter $M_{\mathrm{DM}}(r)$ within a volume of a radius $r$ can be
described by
\begin{equation}
M_{\mathrm{DM}}(r) = 4\pi\rho_0 r_{\mathrm{C}}^2
  \left[r - r_{\mathrm{C}}\arctan(r/r_{\mathrm{C}})\right].
\tag{4}
\end{equation}

\medskip
\textbf{Part 1 The mass of dark matter and rotation curves of the galaxy}

\begin{description}
\item[(D2.1)] In laboratories on Earth, H$\alpha$ has an emitted wavelength
  $\lambda_{\mathrm{emit}}$ of 656.281\,nm. Compute the observed rotational
  velocities of the galaxy $V_{\mathrm{obs}}$ and the rotational velocities
  due to the dark matter $V_{\mathrm{DM}}$ at distance $r$ in units of
  km\,s$^{-1}$.

  For the different values of $r$ given in the table, compute the dynamical
  mass $M(r)$ and the mass of dark matter $M_{\mathrm{DM}}(r)$ in solar
  masses. \hfill[21]
\item[(D2.2)] Create rotation curves of the galaxy on graph paper by plotting
  the points of $V_{\mathrm{obs}}$, $V_{\mathrm{DM}}$, $V_{\mathrm{gas}}$,
  $V_*$ versus the radius $r$ and draw smooth curves through the points (mark
  your graph as ``D2.2'').

  Order the contribution of the different components to the observed velocity
  in descending order. \hfill[16]
\end{description}

\medskip
\textbf{Part 2 Dark matter distribution}

\begin{description}
\item[(D2.3)] Take a data point at small $r$ and large $r$ to estimate
  $\rho_0$ and $r_{\mathrm{C}}$. Note that for large values of $x$,
  $\arctan(x) \approx \pi/2$ and at small $x$,
  $\arctan(x) \approx x - x^3/3$. \hfill[7]
\item[(D2.4)] By comparing Equation (4) to a linear function, the central
  density $\rho_0$ could also be found by a linear fit. Plot an appropriate
  graph so that a linear fit can be used to find another value of $\rho_0$.
  Evaluate $\rho_0$ in units of $M_\odot$\,kpc$^{-3}$. (Mark your graph as
  ``D2.4''). If you cannot find the value of $r_{\mathrm{C}}$ from the
  previous part, use $r_{\mathrm{C}} = 3.2$\,kpc as an estimate for this
  part. \hfill[19]
\item[(D2.5)] Compute logarithmic values of the dark matter density,
  $\ln[\rho_{\mathrm{DM}}(r)]$, and plot the distribution of dark matter in
  the galaxy as a function of radius $r$ on graph paper. (Mark your graph as
  ``D2.5''). \hfill[12]
\end{description}
